\documentclass[11pt]{article}

\usepackage[T1]{fontenc}
\usepackage{amsmath,amssymb,amsthm}
\usepackage{booktabs}
\usepackage[margin=1in]{geometry}
\usepackage{hyperref}
\hypersetup{pdftitle={An Exact Hall-Universality Census Through Six Vertices},pdfauthor={Xiaojian Zeng},pdfsubject={Exact computational graph census},pdfkeywords={Hall condition, graphic list coloring, graph census, exact verification}}
\usepackage{microtype}

\newtheorem{theorem}{Theorem}
\newtheorem{lemma}[theorem]{Lemma}
\newcommand{\cN}{\mathcal N}
\newcommand{\cA}{\mathcal A}

\title{An Exact Hall-Universality Census Through Six Vertices}
\author{Xiaojian Zeng\\
\small PBC School of Finance, Tsinghua University, Beijing 100190, China\\
\small \texttt{xiaojianzeng@foxmail.com}}
\date{}

\begin{document}
\maketitle

\begin{abstract}
For finite simple graphs $G$ and $H$ on the same labelled vertex set, give
vertex $v$ the graphic list $N_G(v)$.  A graph $G$ is Hall-universal when the
Hilton--Johnson Hall condition is sufficient for a proper coloring of every
$H$ from these lists.  We give an exact exhaustive census for $|V(G)|\leq 6$:
every one of the $1,1,2,4,11,34,156$ isomorphism classes in orders zero through
six is Hall-universal.  The certificate stores the Hall-feasibility and
colorability bit of every labelled $H$ for every canonical $G$.  An independent
verifier reconstructs the complete domain of $G$-orbits, tests all vertex
subsets by brute force, and checks colorability directly from admissible
assignments.  Of the 209 positive classes, 199 follow from families identified
in the source paper; the first remaining class in $(|V|,\texttt{graph6})$ order
is the five-vertex bull graph.  This finite theorem does not resolve the
unrestricted characterization problem.
\end{abstract}

\section{Statement and context}

Let $V=\{0,\ldots,n-1\}$, and let $G,H$ be simple undirected graphs on $V$.
A \emph{$G$-coloring} of $H$ is a map $\varphi:V\to V$ such that
\[
  \varphi(v)\in N_G(v)\quad(v\in V),
  \qquad
  \varphi(u)\ne\varphi(v)\quad(uv\in E(H)).
\]
Thus the list attached to $v$ is the open neighborhood of $v$ in $G$.
For $U\subseteq V$ and $c\in V$, put
\[
  a_{H,G}(U,c)=\alpha\bigl(H[U\cap N_G(c)]\bigr),
\]
where $\alpha$ is the independence number and $\alpha(\varnothing)=0$.
Here the vertices eligible for color $c$ are exactly $N_G(c)$, because $G$ is
undirected.
The Hilton--Johnson Hall condition is
\begin{equation}\label{eq:hall}
  |U|\leq \sum_{c\in V}a_{H,G}(U,c)
  \qquad\text{for every }U\subseteq V.
\end{equation}
Its general list-coloring form is necessary for a proper list coloring
\cite{HiltonJohnson1990}; it extends the classical systems-of-representatives
condition of Hall \cite{Hall1935}.  In the graphic-list setting above, Chalise
defined $G$ to be \emph{Hall-universal} if \eqref{eq:hall} implies that $H$ has
a $G$-coloring for every $H$ on the common vertex set
\cite{Chalise2026}.

Chalise proved Hall-universality for several families, including forests,
appropriate complete multipartite graphs, grids, and graphs with a spanning
$\{1,2\}$-factor, and asked for the unrestricted characterization
\cite{Chalise2026}.  The result here concerns only the bounded census requested
through order six.

\begin{theorem}\label{thm:main}
For every integer $n$ with $0\leq n\leq6$, every finite simple graph $G$ on
$n$ vertices is Hall-universal.  Up to isomorphism, the complete positive
census in orders $0,1,\ldots,6$ has respectively
\[
  1,\ 1,\ 2,\ 4,\ 11,\ 34,\ 156
\]
classes, and there are no negative classes.
\end{theorem}

We include the null graph at $n=0$: its unique coloring is empty and its Hall
condition is vacuous.  Under a convention excluding the null graph, only this
one trivial class is removed.

Exact finite graph certification using canonical \texttt{graph6}
representatives and independently checkable per-class data has also been used,
for example, in an eight-vertex whiskered-graph problem
\cite{Zeng2026SSRN7385138}.  Here the objects being certified are different:
for each canonical $G$ we retain two predicates for every labelled $H$, not a
single certificate per isomorphism class.  This comparison is methodological;
it supplies no mathematical input to Theorem~\ref{thm:main}.  The literature
assessment is correspondingly bounded: no claim of priority is made.

\section{Finite certificate}

Order the $\binom n2$ possible edges by increasing larger endpoint and then
increasing smaller endpoint, and identify every labelled $H$ with its edge
mask in $\{0,\ldots,2^{\binom n2}-1\}$.  One canonical \texttt{graph6}
representative is retained for each isomorphism class of $G$.

\subsection{Hall predicate}

For a fixed $H$, the main enumerator obtains all independence numbers from the
integer recurrence
\begin{equation}\label{eq:alpha}
 \alpha_H(S)=\max\bigl\{\alpha_H(S\setminus\{v\}),
 1+\alpha_H((S\setminus\{v\})\setminus N_H(v))\bigr\}.
\end{equation}
It tests \eqref{eq:hall} on every nonempty $U$ for which $H[U]$ is connected.

\begin{lemma}\label{lem:connected}
Condition \eqref{eq:hall} holds for every $U\subseteq V$ if and only if it
holds for every nonempty $U$ such that $H[U]$ is connected.
\end{lemma}

\begin{proof}
Only the reverse implication needs proof.  If the components of $H[U]$ have
vertex sets $U_1,\ldots,U_t$, then for every color $c$ the induced graphs
$H[U_i\cap N_G(c)]$ are mutually anticomplete.  Independence number is
therefore additive over the components.  Summing the inequalities for the
$U_i$ gives the inequality for $U$.  The empty set is immediate.
\end{proof}

This also explains why induced subgraphs suffice in the equivalent formulation
that quantifies over all subgraphs: deleting edges cannot decrease any
independence number.

\subsection{Colorability predicate}

For a fixed $G$, define the finite set of admissible assignments
\[
  \cA_G=\{\varphi:V\to V:\varphi(v)\in N_G(v)\text{ for every }v\}.
\]
For $\varphi\in\cA_G$, its conflict edge set is
\[
  Q_\varphi=\{uv\in\tbinom V2:\varphi(u)=\varphi(v)\}.
\]

\begin{lemma}\label{lem:conflict}
A labelled graph $H$ has a $G$-coloring if and only if
$E(H)\cap Q_\varphi=\varnothing$ for some $\varphi\in\cA_G$.
\end{lemma}

\begin{proof}
Membership in $\cA_G$ is exactly the list constraint.  The displayed
disjointness says exactly that no edge of $H$ has equally colored endpoints.
\end{proof}

The main enumerator visits all at most $n^n$ maps $V\to V$, retains the
admissible ones, and marks the downward closure of their allowable edge masks.
Thus no satisfiability or floating-point oracle enters either predicate.

\subsection{Independent verification}

The independent verifier deliberately avoids the decisive reductions used by
the main enumeration.  It
\begin{enumerate}
  \item enumerates all $2^{\binom n2}$ labelled graphs $G$ and applies all
        vertex permutations, recovering exactly the recorded canonical orbit
        set and hence deriving the seven class counts;
  \item computes each independence number by testing every submask for
        independence, rather than using \eqref{eq:alpha};
  \item checks \eqref{eq:hall} for every nonempty $U$, including disconnected
        sets; and
  \item regenerates the distinct $Q_\varphi$ and tests every $H$ directly
        against them, rather than using a downward closure.
\end{enumerate}
Every operation is on bounded integers and Boolean arrays.  The representative
generation and \texttt{graph6} decoding use NetworkX
\cite{HagbergSchultSwart2008}; completeness is independently established by
the full labelled-$G$ orbit computation rather than assumed from the software's
atlas counts.

For every canonical $G$, the certificate stores two little-endian packed bit
arrays indexed by increasing $H$ edge mask: one for \eqref{eq:hall} and one for
colorability.  The verifier found the arrays identical in every row.  It also
checked row keys, unused padding bits, row bindings, the aggregate binding, and
the complete set of canonical representatives.  The supplementary manifest
records the integrity bindings for the data and verifier inputs.

The complete numerical summary is Table~\ref{tab:census}.  Its ``pairs'' column
counts canonical-$G$/labelled-$H$ pairs, not labelled choices of both graphs.

\begin{table}[ht]
\centering
\small
\begin{tabular}{rrrrrrr}
\toprule
$n$ & $G$ classes & $H$ per class & pairs & Hall & colorable & source-covered $G$\\
\midrule
0 & 1   & 1      & 1         & 1         & 1         & 1\\
1 & 1   & 1      & 1         & 0         & 0         & 1\\
2 & 2   & 2      & 4         & 2         & 2         & 2\\
3 & 4   & 8      & 32        & 12        & 12        & 4\\
4 & 11  & 64     & 704       & 392       & 392       & 11\\
5 & 34  & 1,024  & 34,816    & 19,984    & 19,984    & 31\\
6 & 156 & 32,768 & 5,111,808 & 3,597,344 & 3,597,344 & 149\\
\bottomrule
\end{tabular}
\caption{Exact census.  ``Hall'' and ``colorable'' are counts of positive bits
over canonical-$G$/labelled-$H$ pairs.}
\label{tab:census}
\end{table}

\begin{proof}[Proof of Theorem~\ref{thm:main}]
Simultaneously relabelling the vertices and the colors preserves both
\eqref{eq:hall} and the existence of a $G$-coloring, so one representative of
each $G$-isomorphism class suffices.  The verifier's labelled-orbit computation
shows that the certificate contains every such class for $0\leq n\leq6$.
For each representative it exhausts every labelled $H$.  Lemmas
\ref{lem:connected} and \ref{lem:conflict} justify the two exact predicates,
and the independent verifier recomputes both predicates by the distinct routes
described above.  In every row and at every $H$ mask, a positive Hall bit has a
positive colorability bit.  Hence every represented $G$ is Hall-universal.
The derived orbit counts give the asserted census.
\end{proof}

If $G$ has an isolated vertex $v$, there is no admissible coloring, but Hall
already fails at $U=\{v\}$.  Such classes are therefore positive vacuously;
in particular, the zero Hall and colorability counts at $n=1$ do not hide an
exception.

\section{Comparison with the source families}

A separate exact classifier applies the sufficient families and observations
in the source paper: isolated-vertex vacuity, spanning $\{1,2\}$-factors,
star attachment, forests, the complete multipartite criterion, grids, and the
disjoint-union observation \cite{Chalise2026}.  It classifies 199 of the 209
positive classes.  The remaining ten \texttt{graph6} representatives, in
$(n,\texttt{graph6})$ order, are
\begin{center}
\small\ttfamily
DBk, DBw, DB\{, E?NW, E?No, E?Nw, E?]o, E?]w, E?\char94 o, E?\char94 w.
\end{center}
The first is \texttt{DBk}, edge mask $880$, the bull graph.  In the certificate's
zero-based labels its triangle has edges $23,24,34$ and its pendant edges are
$13,04$.  ``First'' here refers only to the stated reproducible ordering, not
to a priority assertion.

The ten residual classes show why the census is not merely a count of instances
covered by the cited sufficient families.  They do not, however, imply that
those families are exhaustive beyond six vertices, and the unrestricted
characterization remains outside the scope of this theorem.

\section{Reproducibility and assistance disclosure}

The supplementary files contain the canonical representative table, the main
enumerator, the packed per-$H$ census, the source-family classifier, the row
manifest, the independent verifier, and its machine-readable report.  From the
project root, the complete verification is run by
\begin{verbatim}
python3 evidence/verify_census.py evidence/census_exact.tsv \
  evidence/census_manifest.json evidence/representatives.tsv
\end{verbatim}
The original enumeration used Python~3, NetworkX~3.4.2, NumPy~2.2.6, and
Numba~0.67.0.  All proof-essential predicates use exact integer or Boolean
arithmetic.

OpenAI Codex was used during the research workflow to organize the search,
implement and run the exact enumerators, construct a separately routed
verifier, screen the cited sources, and assist in drafting this manuscript.
The full verifier was rerun separately, together with checks of the
quantifiers, bit ordering, orbit completeness, source-family classifier, and
edge cases.  These checks are computational evidence for the explicitly
described finite theorem and do not extend its scope.

\bibliographystyle{plain}
\bibliography{references}

\end{document}
